%!TEX root = ../main.tex
% GLOSARIO — términos conceptuales relevantes de MarketLens.
% Definición ≠ aparición automática: usar \gls{clave} o \glsadd{clave}.

\newglossaryentry{marketplace}{
    name={marketplace},
    description={Plataforma digital que facilita la publicación y el intercambio de bienes o servicios entre compradores y vendedores}
}

\newglossaryentry{prototipo}{
    name={prototipo},
    description={Versión funcional o experimental de un sistema utilizada para comprobar decisiones de diseño y comportamiento en un alcance delimitado; no presupone aptitud para producción}
}

\newglossaryentry{g:mvp}{
    name={producto mínimo viable},
    description={Versión de un producto que reúne las funciones mínimas necesarias para ofrecer una propuesta de valor y aprender de su utilización real}
}

\newglossaryentry{stakeholder}{
    name={parte interesada},
    description={Persona, grupo u organización que puede afectar al proyecto o verse afectada por sus resultados o decisiones}
}

\newglossaryentry{indicador-riesgo}{
    name={indicador de riesgo},
    description={Atributo observable o derivado que se utiliza como señal para estimar la necesidad de revisar una publicación; por sí solo no constituye prueba de fraude}
}

\newglossaryentry{indice-riesgo}{
    name={índice de riesgo},
    description={Puntuación calculada por MarketLens mediante reglas y ponderaciones; representa una evaluación orientativa de factores presentes, no una probabilidad estadística de fraude}
}

\newglossaryentry{cobertura-informacion}{
    name={cobertura de información},
    description={Proporción o recuento de grupos de datos para los que se dispone de información suficiente durante una evaluación; permite distinguir evidencia ausente de señales negativas observadas}
}

\newglossaryentry{referencia-provisional}{
    name={referencia provisional de precio},
    description={Publicación candidata que supera determinados filtros de comparabilidad, sin garantizar equivalencia comercial perfecta ni suficiencia estadística para estimar precios de mercado}
}

\newglossaryentry{normalizacion}{
    name={normalización de productos},
    description={Tratamiento de títulos y atributos para reducir variaciones de escritura y facilitar la comparación entre publicaciones potencialmente equivalentes}
}

\newglossaryentry{similitud-coseno}{
    name={similitud coseno},
    description={Medida de similitud entre representaciones vectoriales que considera el coseno del ángulo formado por los vectores; utilizada para comparar títulos de publicaciones}
}

\newglossaryentry{falso-positivo}{
    name={falso positivo},
    description={Caso legítimo clasificado incorrectamente como sospechoso o riesgoso por un mecanismo de evaluación}
}

\newglossaryentry{falso-negativo}{
    name={falso negativo},
    description={Caso problemático que un mecanismo de evaluación no identifica como riesgoso}
}

\newglossaryentry{exposicion}{
    name={exposición al riesgo},
    description={Producto de una valoración ordinal de probabilidad por una valoración ordinal de impacto; sirve para priorizar la gestión de riesgos del proyecto, sin representar una probabilidad medida}
}

\newglossaryentry{g:van}{
    name={valor actual neto},
    description={Diferencia entre el valor presente de los flujos de caja proyectados y la inversión inicial, obtenida mediante una tasa de descuento asumida; su interpretación depende de la validez de los supuestos}
}
